Systems that generate GPU kernels with language models often time several
candidates and promote the fastest. This couples search to evaluation. Timing
noise and evaluator state can then be mistaken for performance. We present
\system{}, an evaluation protocol that freezes candidates before timing and
confirms the selected winner in fresh processes. It combines persistent modules,
contract checks, randomized paired CUDA events, formal controls, and checksummed
artifacts. We apply the protocol to 48 KernelBench L1 tasks selected without
performance data. Of 141 evaluated Gemini candidates, 27 pass every gate and
cover 10 tasks. No selected task winner clears a prespecified 2\% margin over
eager execution in screening or confirmation. One candidate confirms a
2.001$\times$ gain over \code{torch.compile} while remaining below eager. A
deterministic control with 80 easy candidates yields no false promotions. In a
calibrated stress test near parity, screening promotes three of four task winners
at $K=8$, but only two confirm. The interval across four tasks includes zero, so
this is evidence for a mechanism rather than a population estimate. A lifecycle
ablation finds 1.053$\times$ median host latency inflation without a corresponding
CUDA event change. Independent confirmation can therefore change the claim
supported by a kernel search.
